% !TEX root = ../main.tex

%----------------------------------------------------------------------------------------
% CHAPTER 8 - ANÀLISI I DISSENY DEL SISTEMA
%----------------------------------------------------------------------------------------

\chapter{Anàlisi i disseny del sistema}
\label{Ch:AnalysisDesign}

%----------------------------------------------------------------------------------------

\section{Anàlisi del sistema}

El sistema de visualització de dades hídriques de Catalunya neix de la necessitat de proporcionar accés públic, intuïtiu i actualitzat a les dades d'embassaments i precipitació que la Generalitat de Catalunya publica a través de les seves APIs Socrata. L'anàlisi detallada de les necessitats del sistema revela diversos aspectes crítics:

\begin{itemize}
    \item \textbf{Accessibilitat de les dades}: Les dades hídriques estan disponibles però disperses en múltiples endpoints. Els usuaris han de navegar per diferents portals i formats per obtenir una visió completa. El sistema ha de consolidar aquesta informació en una única interfície coherent.
    
    \item \textbf{Actualització contínua}: Els nivells dels embassaments s'actualitzen cada 15 minuts i les dades de precipitació cada hora. El sistema ha de garantir que la informació visualitzada sigui sempre la més recent possible sense requerir intervenció manual.
    
    \item \textbf{Correlació de dades}: Per entendre l'evolució dels recursos hídrics és essencial poder correlacionar els nivells d'embassaments amb les dades de precipitació de les estacions meteorològiques més properes.
    
    \item \textbf{Detecció de situacions crítiques}: En períodes de sequera, és crucial poder identificar ràpidament quins embassaments es troben en situacions de risc mitjançant indicadors visuals clars.
    
    \item \textbf{Visualització geogràfica}: La ubicació geogràfica dels embassaments i estacions meteorològiques és rellevant per entendre la distribució territorial dels recursos hídrics i les precipitacions.
\end{itemize}

Aquestes necessitats han portat al disseny d'un sistema amb quatre panells de visualització diferenciats, una canonada ETL automatitzada i una arquitectura que separa clarament el processament de dades de la seva presentació.

%----------------------------------------------------------------------------------------

\subsection{Tipus d'usuari i rols}

El sistema està dissenyat per atendre tres perfils d'usuari principals, cadascun amb necessitats i nivells de coneixement tècnic diferents:

\begin{enumerate}
    \item \textbf{Ciutadà interessat}: Usuari no tècnic que vol consultar l'estat actual dels embassaments de Catalunya, entendre si hi ha situacions de sequera i visualitzar tendències bàsiques. Necessita interfícies intuïtives, indicadors visuals clars (com els semàfors d'alerta) i explicacions accessibles. Aquest perfil representa la majoria d'usuaris potencials.
    
    \item \textbf{Professional o investigador}: Usuari amb coneixements tècnics que necessita accedir a dades detallades, correlacionar variables (precipitació vs nivells), analitzar tendències temporals i possiblement exportar dades per a anàlisis externes. Valorarà la precisió de les dades, la possibilitat de filtrar per períodes i estacions, i la qualitat de les visualitzacions.
    
    \item \textbf{Administrador del sistema}: Responsable de mantenir la canonada ETL, verificar que les actualitzacions automàtiques funcionen correctament, i assegurar la disponibilitat del servei. Necessita logs detallats, mecanismes de monitorització i la capacitat de reexecutar processos en cas d'error.
\end{enumerate}

El disseny del sistema prioritza l'experiència del ciutadà interessat, assegurant que totes les funcionalitats siguin accessibles sense coneixements tècnics previs, mentre que els perfils professionals troben eines d'anàlisi més avançades als panells de correlació i alertes.

%----------------------------------------------------------------------------------------

\subsection{Anàlisi de requisits}

A partir dels requisits definits al Capítol 5, s'ha realitzat una anàlisi detallada de què implica cada requisit i quines dades cal gestionar:

\subsubsection{Requisits funcionals}

\textbf{RF-1 (Extracció automàtica de dades)}: Implica la implementació de clients HTTP que realitzin peticions GET als endpoints de Socrata per als datasets d'embassaments i precipitació. Cal gestionar la paginació de resultats, els límits de taxa (1.000 peticions/hora), i els possibles errors de xarxa. Les dades obtingudes inclouen per a cada embassament: nom, percentatge d'ocupació, volum embassat (hm³), cota (msnm), i data d'actualització. Per a les estacions meteorològiques: codi, nom, precipitació (mm), i data de mesura.

\textbf{RF-2 (Transformació i neteja)}: Les dades brutes de Socrata contenen inconsistències com valors nuls, formats de data heterogenis, i unitats que poden requerir conversió. El transformador ha de normalitzar totes les dates a format ISO 8601, eliminar registres duplicats, gestionar valors que falten (marcant-los explícitament), i estructurar les dades en un format JSON estàndard que el frontend pugui consumir directament.

\textbf{RF-3 (Emmagatzematge en JSON versionat)}: Les dades processades es desen en fitxers JSON dins del repositori Git, amb timestamps que permeten mantenir l'historial. Això implica una estructura de carpetes clara (\texttt{/public/data/}) i noms de fitxers que indiquin el tipus de dada i la data. El versionat amb Git proporciona automàticament traçabilitat de canvis.

\textbf{RF-4 (Quatre panells de visualització)}: Cada panell requereix components React específics:
\begin{itemize}
    \item Panell 1 (Mapa): Integració de Leaflet amb markers per a cada embassament i estació, popups informatius, i KPIs dinàmics.
    \item Panell 2 (Evolució temporal): Gràfics de línies amb Recharts, selectors de data, i funcionalitat de comparació múltiple.
    \item Panell 3 (Correlació): Gràfics de dispersió amb càlcul del coeficient de Pearson i selecció d'estacions.
    \item Panell 4 (Alertes): Taula amb indicadors semàfor basats en thresholds predefinits i fletxes de tendència.
\end{itemize}

\textbf{RF-5 (Detecció de sequera)}: Requereix implementar la lògica de classificació basada en percentatge d'ocupació: crític (<20\%), avís (20-50\%), normal (50-75\%), òptim (>75\%). A més, s'ha afegit un algoritme de regressió lineal sobre les últimes 6 mesures per detectar tendències de millora o deteriorament.

\textbf{RF-6 (Automatització amb GitHub Actions)}: Implica la configuració d'un workflow YAML amb programació cron (0 2 * * *) que executi l'script ETL diàriament. Cal gestionar el checkout del repositori, la instal·lació de dependències, l'execució de l'script, i el commit/push dels nous fitxers JSON.

\subsubsection{Requisits no funcionals}

\textbf{RNF-1 (Rendiment)}: L'ETL ha de processar >200.000 registres en total en <5 minuts (utilitzant peticions paginades de 50.000 registres). Això implica optimitzar les peticions HTTP (paral·lelisme quan sigui possible), minimitzar el processament redundant, i utilitzar operacions de fitxer eficients. El frontend ha de carregar en <3 segons, requerint optimitzacions com la minificació de bundles, la compressió d'assets, i el carregament diferit de components no crítics.

\textbf{RNF-2 (Usabilitat)}: La interfície ha de ser intuïtiva per a usuaris no tècnics. Això implica navegació clara amb etiquetes descriptives, elements interactius amb feedback visual immediat, i disseny responsive que funcioni en dispositius mòbils i d'escriptori.

\textbf{RNF-3 (Seguretat)}: Totes les connexions a les APIs són HTTPS. No s'emmagatzemen credencials al repositori (les APIs són públiques). La configuració externalitzada amb \texttt{.env} permet gestionar paràmetres sensibles si calgués en el futur.

\textbf{RNF-4 (Portabilitat)}: El frontend és una aplicació estàtica que es pot desplegar a qualsevol servidor web (GitHub Pages, Vercel, Netlify). L'ETL és un script Node.js que pot executar-se en qualsevol entorn amb Node.js 18+.

\textbf{RNF-5 (Mantenibilitat)}: Codi modular amb separació clara de responsabilitats, convencions de nomenclatura consistents, comentaris explicatius, i documentació tècnica actualitzada.

\textbf{RNF-6 (Escalabilitat)}: L'arquitectura permet afegir nous tipus de dades o augmentar el volum sense redesign complet. La separació ETL/frontend i l'ús de fitxers JSON com a intermediari faciliten aquesta escalabilitat.

%----------------------------------------------------------------------------------------

\subsection{Casos d'ús}

Els principals casos d'ús del sistema es resumeixen a la taula següent:

\begin{table}[h]
\centering
\caption{Resum de casos d'ús del sistema}
\begin{tabular}{|l|l|p{6cm}|}
\hline
\textbf{ID} & \textbf{Actor} & \textbf{Descripció} \\
\hline
CU-1 & Ciutadà, Professional & Visualitzar l'estat actual dels embassaments al mapa interactiu, amb markers que mostren percentatge d'ocupació i informació detallada al clicar. \\
\hline
CU-2 & Ciutadà, Professional & Consultar l'evolució temporal del nivell d'un embassament amb gràfic de línies, selector de dates i comparació múltiple. \\
\hline
CU-3 & Professional, Investigador & Analitzar la correlació entre precipitació i nivells mitjançant gràfics de dispersió i coeficient de Pearson, amb selecció d'estacions i períodes. \\
\hline
CU-4 & Ciutadà, Professional & Consultar l'estat d'alerta de sequera amb indicadors semàfor (crític, avís, normal, òptim) i fletxes de tendència. \\
\hline
CU-5 & Sistema (GitHub Actions & Executar l'actualització automàtica diària: extreure dades de Socrata, transformar, desa-les en JSON i fer push al repositori. \\
\hline
\end{tabular}
\end{table}

%----------------------------------------------------------------------------------------

\subsection{Model conceptual de dades}

El sistema gestiona principalment tres entitats interrelacionades: embassaments, estacions meteorològiques i mesures. El model conceptual es pot representar mitjançant un diagrama entitat-relació simplificat:

\begin{itemize}
    \item \textbf{Embassament} (\texttt{Reservoir}):
    \begin{itemize}
        \item \texttt{id} (identificador únic)
        \item \texttt{nom} (nom de l'embassament)
        \item \texttt{percentatge\_ocupacio} (\%)
        \item \texttt{volum\_embassat} (hm³)
        \item \texttt{cota} (msnm)
        \item \texttt{capacitat\_total} (hm³)
        \item \texttt{latitud}, \texttt{longitud} (coordenades geogràfiques)
        \item \texttt{ultima\_actualitzacio} (timestamp)
    \end{itemize}
    
    \item \textbf{Estació meteorològica} (\texttt{WeatherStation}):
    \begin{itemize}
        \item \texttt{codi} (identificador únic)
        \item \texttt{nom} (nom de l'estació)
        \item \texttt{latitud}, \texttt{longitud} (coordenades geogràfiques)
        \item \texttt{ultima\_actualitzacio} (timestamp)
    \end{itemize}
    
    \item \textbf{Mesura de precipitació} (\texttt{PrecipitationMeasurement}):
    \begin{itemize}
        \item \texttt{estacio\_codi} (referència a Estació)
        \item \texttt{data\_mesura} (timestamp)
        \item \texttt{precipitacio\_mm} (mm)
    \end{itemize}
    
    \item \textbf{Alerta de sequera} (\texttt{DroughtAlert}):
    \begin{itemize}
        \item \texttt{embassament\_id} (referència a Embassament)
        \item \texttt{data} (timestamp)
        \item \texttt{percentatge\_ocupacio} (\%)
        \item \texttt{nivell\_alerta} (crític, avís, normal, òptim)
        \item \texttt{tendencia} (millora, estable, deteriorament)
    \end{itemize}
\end{itemize}

Les relacions són:
\begin{itemize}
    \item Una \textbf{Estació} pot tenir múltiples \textbf{Mesures de precipitació} (1:N)
    \item Un \textbf{Embassament} pot generar múltiples \textbf{Alertes de sequera} al llarg del temps (1:N)
    \item Les \textbf{Mesures} i els \textbf{Embassaments} es relacionen indirectament a través de la proximitat geogràfica per a l'anàlisi de correlació
\end{itemize}

Aquest model es materialitza en fitxers JSON estructurats que el frontend consumeix directament, sense necessitat d'una base de dades relacional.

%----------------------------------------------------------------------------------------

\section{Disseny del sistema}

La solució proposada per al sistema de visualització de dades hídriques es basa en una arquitectura desacoblada que separa clarament el processament de dades (ETL) de la seva presentació (frontend). Aquesta separació proporciona robustesa, flexibilitat de desplegament i facilitat de manteniment.

El sistema es compon de tres components principals:

\begin{enumerate}
    \item \textbf{Canonada ETL}: Scripts Node.js que s'executen periòdicament per extreure dades de les APIs de Socrata, transformar-les en un format estàndard, i desar-les com a fitxers JSON versionats al repositori Git.
    
    \item \textbf{Repositori de dades}: Fitxers JSON estructurats dins del repositori que actuen com a font de dades per al frontend. El versionat amb Git proporciona historial i traçabilitat automàtics.
    
    \item \textbf{Frontend React}: Aplicació web de pàgina única (SPA) que consumeix els fitxers JSON i els presenta mitjançant quatre panells de visualització interconnectats, amb navegació fluida entre ells.
\end{enumerate}

Aquesta arquitectura té l'avantatge que no requereix un servidor de backend en producció: el frontend és completament estàtic i es pot servir des de qualsevol CDN o servidor web estàtic. L'ETL s'executa en l'entorn de GitHub Actions, aprofitant la infraestructura gratuïta disponible per a repositoris públics.

%----------------------------------------------------------------------------------------

\subsection{Arquitectura del sistema}

L'arquitectura general del sistema segueix un patró de pipeline de dades amb separació clara entre les capes d'extracció, emmagatzematge i presentació:

\begin{enumerate}
    \item \textbf{Capa d'extracció (Extract)}:
    \begin{itemize}
        \item Mòdul \texttt{extractReservoirs}: Realitza peticions HTTP GET a l'endpoint de Socrata per a dades d'embassaments
        \item Mòdul \texttt{extractPrecipitation}: Realitza peticions HTTP GET a l'endpoint de Socrata per a dades de precipitació
        \item Gestió centralitzada d'errors, reintents i logging
    \end{itemize}
    
    \item \textbf{Capa de transformació (Transform)}:
    \begin{itemize}
        \item Normalització de formats de data a ISO 8601
        \item Eliminació de duplicats i registres inconsistents
        \item Càlcul de mètriques derivades (percentatges, tendències)
        \item Classificació d'alertes segons thresholds predefinits
        \item Estructuració en format JSON estàndard
    \end{itemize}
    
    \item \textbf{Capa de càrrega (Load)}:
    \begin{itemize}
        \item Escriptura de fitxers JSON a \texttt{/public/data/}
        \item Generació de fitxers d'índex per a càrrega eficient
        \item Commit i push automàtic dels canvis al repositori
    \end{itemize}
    
    \item \textbf{Capa de presentació (Frontend)}:
    \begin{itemize}
        \item Servei de dades (\texttt{waterDataService}): Wrapper amb axios per carregar fitxers JSON
        \item Hooks personalitzats (\texttt{useWaterData}): Gestió d'estat de càrrega, error i dades
        \item Store global (Zustand): Estat compartit entre components
        \item Components de dashboards: Quatre panells amb visualitzacions específiques
        \item Sistema de navegació: React Router per a transicions entre panells
    \end{itemize}
\end{enumerate}

El flux de dades és unidireccional: les APIs $\rightarrow$ ETL $\rightarrow$ JSON $\rightarrow$ Frontend $\rightarrow$ Usuari. No hi ha comunicació bidireccional ni escriptura des del frontend, simplificant l'arquitectura i reduint els punts de fallada.

%----------------------------------------------------------------------------------------

\subsection{Disseny de les interfícies d'usuari}

La filosofia de disseny de les interfícies segueix els principis de \textbf{disseny responsiu} (responsive design) amb enfocament \textbf{mobile-first}, assegurant que l'experiència sigui òptima en tots els dispositius:

\begin{itemize}
    \item \textbf{Simplicitat i claredat}: Cada panell té un objectiu clar i evita la sobrecàrrega visual. Els elements es presenten de forma jeràrquica, amb la informació més important destacada.
    
    \item \textbf{Consistència visual}: S'utilitza una paleta de colors coherent amb el tema hídric (tons de blau i verd), tipografia uniforme, i espaiament consistent entre elements. Les variables SCSS centralitzades a \texttt{\_variables.scss} garanteixen aquesta consistència.
    
    \item \textbf{Feedback visual immediat}: Els estats de càrrega es mostren amb spinners, els errors amb missatges clars, i les accions de l'usuari amb transicions suaus. Els indicadors semàfor del Panell 4 proporcionen informació instantània sobre l'estat de sequera.
    
    \item \textbf{Accesibilitat}: S'utilitzen etiquetes descriptives, contrast adequat entre text i fons, i elements interactius prou grans per a ser clicables en dispositius tàctils.
    
    \item \textbf{Metodologia BEM}: L'estructuració CSS segueix la metodologia Block-Element-Modifier, facilitant la mantenibilitat i evitant col·lisions de noms de classes. Per exemple: \texttt{dashboard\_\_card--active} indica un element \texttt{card} dins del bloc \texttt{dashboard} amb el modificador \texttt{active}.
    
    \item \textbf{Grid i Flexbox}: S'utilitza CSS Grid per a layouts principals i Flexbox per a alineació de components interns, proporcionant flexibilitat i adaptabilitat a diferents mides de pantalla.
\end{itemize}

El disseny evita captures de pantalla excessives i se centra en proporcionar una experiència fluida on l'usuari pugui navegar entre els quatre panells sense perdre el context dels filtres aplicats.

%----------------------------------------------------------------------------------------

\subsection{Patrons de disseny}

S'han aplicat diversos patrons de disseny per garantir un codi mantenible i escalable:

\begin{table}[h]
\centering
\caption{Patrons de disseny aplicats}
\begin{tabular}{|l|l|p{6cm}|}
\hline
\textbf{Patró} & \textbf{Ubicació} & \textbf{Benefici principal} \\
\hline
Pipeline modular & ETL & Fases independents (Extract/Transform/Load) que es poden provar i modificar de forma aïllada. \\
\hline
Servei (Service Layer) & Frontend & Centralització de l'accés a dades en \texttt{waterDataService.js}, facilitant el canvi de font de dades. \\
\hline
Basat en components & Frontend & Reutilització de components React, proves aïllades i mantenibilitat millorada. \\
\hline
Observador (Zustand) & Frontend & Desacoblement entre components i actualització automàtica de l'interfície. \\
\hline
Decorador (Middleware) & ETL & Afegir logging i gestió d'errors sense modificar la lògica de negoci. \\
\hline
Strategy & ETL & Intercanviar algoritmes de transformació segons el tipus de dada. \\
\hline
\end{tabular}
\end{table}

Aquests patrons treballen conjuntament per crear un sistema robust que pot adaptar-se a futurs requisits sense redissenys importants.
